\documentclass[a4paper,twoside]{article}

\usepackage{morewrites}

\usepackage[svgnames,dvipsnames,x11names]{xcolor}
\usepackage{graphicx}
\usepackage{texsmith-lists}
\usepackage[english]{babel}
\usepackage[colorlinks=true, linkcolor=NavyBlue, urlcolor=NavyBlue, citecolor=NavyBlue]{hyperref}
\usepackage[a4paper]{geometry}
\usepackage{ts-fonts}
\usepackage{booktabs}
\usepackage{tabularx}
\usepackage{amsmath}
\usepackage{float}
\usepackage{seqsplit}
\newcolumntype{Y}{>{\centering\arraybackslash}X}
\newcolumntype{Z}{>{\raggedleft\arraybackslash}X}
\usepackage{ts-callouts}
\usepackage{ts-code}

\usepackage{lastpage}
\usepackage{refcount}
\makeatletter
\def\ps@plain{
\let\@oddhead\@empty
\let\@evenhead\@empty
\def\@oddfoot{
\hfil
\ifnum\getpagerefnumber{LastPage}>1\relax
\thepage
\fi
\hfil}
\let\@evenfoot\@oddfoot
}
\makeatother

\title{Templates}
\author{}\date{}
\begin{document}
\maketitle
TeXSmith templates define the \LaTeX{} skeleton that wraps converted Markdown. Each
template bundles assets, slot definitions, attribute schemas, and build metadata
so you can aim the renderer at anything from articles to slide decks.

Templates are a core feature that allows TeXSmith to be extended for different
use cases. You may use:
\begin{itemize}
\item Custom project templates stored in a \tscodeinline{templates/} folder.
\item Local user templates installed in your \tscodeinline{\textasciitilde{}/.texsmith/templates/} folder.
\item Published templates installed from PyPI.
\item Built-in templates included with TeXSmith.
\end{itemize}
\section{Built-in templates}
TeXSmith includes standard templates for common document types:
\begin{itemize}
\item \tscodeinline{article}: academic-style article layout with title page, abstract, and sections.
\item \tscodeinline{book}: book-style layout with parts, chapters, preface, table of contents, appendices, and optional \tscodeinline{part} base-level override.
\item \tscodeinline{letter}: formal letters with sender/recipient metadata, fold marks, and signature.
\item \tscodeinline{snippet}: standalone \tscodeinline{tikzpicture} frame for screenshots, stickers, and snippets.
\end{itemize}
Invoke them in the CLI with \tscodeinline{-\allowbreak{}tarticle} or \tscodeinline{-\allowbreak{}tletter}. Additional community templates remain available under \tscodeinline{templates/} or as separate PyPI packages.

Take a look at built-in fragments for full configuration options.
\section{Quick start}
\begin{tscode}[lang=bash, engine=pygments]
\begin{Verbatim}[commandchars=\\\{\},breaklines, breakanywhere, commandchars=\\\{\}]
\PY{c+c1}{\PYZsh{} 1. Inspect a built\PYZhy{}in template}
texsmith\PY{+w}{ }\PYZhy{}\PYZhy{}template\PY{+w}{ }article\PY{+w}{ }\PYZhy{}\PYZhy{}template\PYZhy{}info

\PY{c+c1}{\PYZsh{} 2. Render a document with template slots + build}
texsmith\PY{+w}{ }docs/intro.md\PY{+w}{ }\PY{l+s+se}{\PYZbs{}}
\PY{+w}{  }\PYZhy{}\PYZhy{}template\PY{+w}{ }article\PY{+w}{ }\PY{l+s+se}{\PYZbs{}}
\PY{+w}{  }\PYZhy{}\PYZhy{}slot\PY{+w}{ }mainmatter:@document\PY{+w}{ }\PY{l+s+se}{\PYZbs{}}
\PY{+w}{  }\PYZhy{}\PYZhy{}output\PYZhy{}dir\PY{+w}{ }build/article\PY{+w}{ }\PY{l+s+se}{\PYZbs{}}
\PY{+w}{  }\PYZhy{}\PYZhy{}build
\end{Verbatim}
\end{tscode}
Need a blank template skeleton? Use the scaffold flag to copy any template into
your working tree:
\begin{tscode}[lang=bash, engine=pygments]
\begin{Verbatim}[commandchars=\\\{\},breaklines, breakanywhere, commandchars=\\\{\}]
texsmith\PY{+w}{ }\PYZhy{}\PYZhy{}template\PY{+w}{ }article\PY{+w}{ }\PYZhy{}\PYZhy{}template\PYZhy{}scaffold\PY{+w}{ }my\PYZhy{}template
\PY{n+nb}{cd}\PY{+w}{ }my\PYZhy{}template
uv\PY{+w}{ }run\PY{+w}{ }hatch\PY{+w}{ }build\PY{+w}{  }\PY{c+c1}{\PYZsh{} optional packaging smoke test}
\end{Verbatim}
\end{tscode}
Publish by pointing \tscodeinline{pyproject.toml} to the \tscodeinline{template/} package, then \tscodeinline{uv publish} or \tscodeinline{twine upload dist/*}.
For in-depth patterns (overrides, slots, metadata), see the Template Cookbook.
\section{Write your own TeXSmith templates}
TeXSmith uses Jinja2 templates to assemble \LaTeX{} documents from converted
Markdown fragments. You can create your own templates to control the layout,
styling, and structure of the final output.

You may start by copying an existing template package and modifying it to suit
your needs. Templates can either be published on PyPI for easy distribution or kept locally for personal use. By default, TeXSmith will look for templates installed in the current Python environment or in the current working directory under a \tscodeinline{templates/} folder.
\subsection{Structure}
\begin{tscode}[lang=text, stretch=0.5, engine=pygments]
\begin{Verbatim}[commandchars=\\\{\},breaklines, breakanywhere, commandchars=\\\{\}]
├── README.md
├── \PYZus{}\PYZus{}init\PYZus{}\PYZus{}.py          \PYZsh{} optional: a WrappableTemplate subclass, custom normalisers
├── pyproject.toml       \PYZsh{} only to publish the template as a package
└── template
    ├── assets
    │   └── latexmkrc
    ├── manifest.toml
    ├── template.tex     \PYZsh{} the LaTeX entrypoint
    ├── template.typ     \PYZsh{} optional: the Typst scaffolding, [typst.template]
    ├── mermaid\PYZhy{}config.json  \PYZsh{} optional
    └── .sty, .cls...
\end{Verbatim}
\end{tscode}
\tscodeinline{texsmith -\allowbreak{}-\allowbreak{}template article -\allowbreak{}-\allowbreak{}template-\allowbreak{}scaffold DEST} writes exactly this
shape for a built-in, which is the quickest way to see a real one.
\subsection{TOML manifest}
The \tscodeinline{manifest.toml} file describes the template and its metadata.
\begin{tscode}[lang=toml, engine=pygments]
\begin{Verbatim}[commandchars=\\\{\},breaklines, breakanywhere, commandchars=\\\{\}]
\PY{k}{[}\PY{k}{compat}\PY{k}{]}
\PY{c+c1}{\PYZsh{} Specify the compatible TeXSmith versions. Do not forget}
\PY{c+c1}{\PYZsh{} to set a maximum version to avoid future incompatibilities.}
\PY{n}{texsmith}\PY{+w}{ }\PY{o}{=}\PY{+w}{ }\PY{l+s+s2}{\PYZdq{}}\PY{l+s+s2}{\PYZgt{}=0.1,\PYZlt{}1.0}\PY{l+s+s2}{\PYZdq{}}

\PY{k}{[}\PY{k}{latex}\PY{k}{.}\PY{k}{template}\PY{k}{]}
\PY{n}{name}\PY{+w}{ }\PY{o}{=}\PY{+w}{ }\PY{l+s+s2}{\PYZdq{}}\PY{l+s+s2}{acme}\PY{l+s+s2}{\PYZdq{}}
\PY{n}{version}\PY{+w}{ }\PY{o}{=}\PY{+w}{ }\PY{l+s+s2}{\PYZdq{}}\PY{l+s+s2}{0.1.0}\PY{l+s+s2}{\PYZdq{}}
\PY{n}{entrypoint}\PY{+w}{ }\PY{o}{=}\PY{+w}{ }\PY{l+s+s2}{\PYZdq{}}\PY{l+s+s2}{template/template.tex}\PY{l+s+s2}{\PYZdq{}}
\PY{c+c1}{\PYZsh{} IR passes this template needs, resolved with the same import machinery as}
\PY{c+c1}{\PYZsh{} attribute normalisers and applied only while this template renders.}
\PY{n}{passes}\PY{+w}{ }\PY{o}{=}\PY{+w}{ }\PY{p}{[}\PY{l+s+s2}{\PYZdq{}}\PY{l+s+s2}{acme\PYZus{}pkg.questions:run}\PY{l+s+s2}{\PYZdq{}}\PY{p}{]}
\PY{c+c1}{\PYZsh{} Useful for generic docker images or CI pipelines, where}
\PY{c+c1}{\PYZsh{} we only install the required packages.}
\PY{n}{texlive\PYZus{}year}\PY{+w}{ }\PY{o}{=}\PY{+w}{ }\PY{l+m+mi}{2023}
\PY{n}{tlmgr\PYZus{}packages}\PY{+w}{ }\PY{o}{=}\PY{+w}{ }\PY{p}{[}
\PY{+w}{    }\PY{l+s+s2}{\PYZdq{}}\PY{l+s+s2}{babel}\PY{l+s+s2}{\PYZdq{}}\PY{p}{,}
\PY{+w}{    }\PY{l+s+s2}{\PYZdq{}}\PY{l+s+s2}{geometry}\PY{l+s+s2}{\PYZdq{}}\PY{p}{,}
\PY{+w}{    }\PY{l+s+s2}{\PYZdq{}}\PY{l+s+s2}{hyperref}\PY{l+s+s2}{\PYZdq{}}\PY{p}{,}
\PY{+w}{    }\PY{l+s+s2}{\PYZdq{}}\PY{l+s+s2}{microtype}\PY{l+s+s2}{\PYZdq{}}\PY{p}{,}
\PY{+w}{    }\PY{l+s+s2}{\PYZdq{}}\PY{l+s+s2}{fontspec}\PY{l+s+s2}{\PYZdq{}}\PY{p}{,}
\PY{+w}{    }\PY{l+s+s2}{\PYZdq{}}\PY{l+s+s2}{biblatex}\PY{l+s+s2}{\PYZdq{}}\PY{p}{,}
\PY{p}{]}

\PY{c+c1}{\PYZsh{} Attributes:}
\PY{c+c1}{\PYZsh{} \PYZhy{}\PYZhy{}\PYZhy{}\PYZhy{}\PYZhy{}\PYZhy{}\PYZhy{}\PYZhy{}\PYZhy{}\PYZhy{}\PYZhy{}}
\PY{c+c1}{\PYZsh{} Attributes are declared as nested tables describing the type, default}
\PY{c+c1}{\PYZsh{} value, and optional normalisation rules for each entry. Types are inferred}
\PY{c+c1}{\PYZsh{} from the default when omitted.}
\PY{k}{[}\PY{k}{latex}\PY{k}{.}\PY{k}{template}\PY{k}{.}\PY{k}{attributes}\PY{k}{.}\PY{k}{base\PYZus{}level}\PY{k}{]}
\PY{n}{default}\PY{+w}{ }\PY{o}{=}\PY{+w}{ }\PY{l+m+mi}{1}
\PY{n}{type}\PY{+w}{ }\PY{o}{=}\PY{+w}{ }\PY{l+s+s2}{\PYZdq{}}\PY{l+s+s2}{integer}\PY{l+s+s2}{\PYZdq{}}

\PY{k}{[}\PY{k}{latex}\PY{k}{.}\PY{k}{template}\PY{k}{.}\PY{k}{attributes}\PY{k}{.}\PY{k}{title}\PY{k}{]}
\PY{n}{default}\PY{+w}{ }\PY{o}{=}\PY{+w}{ }\PY{l+s+s2}{\PYZdq{}}\PY{l+s+s2}{Title}\PY{l+s+s2}{\PYZdq{}}
\PY{n}{type}\PY{+w}{ }\PY{o}{=}\PY{+w}{ }\PY{l+s+s2}{\PYZdq{}}\PY{l+s+s2}{string}\PY{l+s+s2}{\PYZdq{}}
\PY{n}{escape}\PY{+w}{ }\PY{o}{=}\PY{+w}{ }\PY{l+s+s2}{\PYZdq{}}\PY{l+s+s2}{latex}\PY{l+s+s2}{\PYZdq{}}
\PY{n}{allow\PYZus{}empty}\PY{+w}{ }\PY{o}{=}\PY{+w}{ }\PY{k+kc}{false}

\PY{k}{[}\PY{k}{latex}\PY{k}{.}\PY{k}{template}\PY{k}{.}\PY{k}{attributes}\PY{k}{.}\PY{k}{subtitle}\PY{k}{]}
\PY{n}{default}\PY{+w}{ }\PY{o}{=}\PY{+w}{ }\PY{l+s+s2}{\PYZdq{}}\PY{l+s+s2}{\PYZdq{}}
\PY{n}{type}\PY{+w}{ }\PY{o}{=}\PY{+w}{ }\PY{l+s+s2}{\PYZdq{}}\PY{l+s+s2}{string}\PY{l+s+s2}{\PYZdq{}}
\PY{n}{escape}\PY{+w}{ }\PY{o}{=}\PY{+w}{ }\PY{l+s+s2}{\PYZdq{}}\PY{l+s+s2}{latex}\PY{l+s+s2}{\PYZdq{}}

\PY{k}{[}\PY{k}{latex}\PY{k}{.}\PY{k}{template}\PY{k}{.}\PY{k}{attributes}\PY{k}{.}\PY{k}{author}\PY{k}{]}
\PY{n}{default}\PY{+w}{ }\PY{o}{=}\PY{+w}{ }\PY{l+s+s2}{\PYZdq{}}\PY{l+s+s2}{John Doe}\PY{l+s+s2}{\PYZdq{}}
\PY{n}{type}\PY{+w}{ }\PY{o}{=}\PY{+w}{ }\PY{l+s+s2}{\PYZdq{}}\PY{l+s+s2}{string}\PY{l+s+s2}{\PYZdq{}}
\PY{n}{escape}\PY{+w}{ }\PY{o}{=}\PY{+w}{ }\PY{l+s+s2}{\PYZdq{}}\PY{l+s+s2}{latex}\PY{l+s+s2}{\PYZdq{}}

\PY{k}{[}\PY{k}{latex}\PY{k}{.}\PY{k}{template}\PY{k}{.}\PY{k}{attributes}\PY{k}{.}\PY{k}{date}\PY{k}{]}
\PY{n}{default}\PY{+w}{ }\PY{o}{=}\PY{+w}{ }\PY{l+s+s2}{\PYZdq{}}\PY{l+s+se}{\PYZbs{}\PYZbs{}}\PY{l+s+s2}{today}\PY{l+s+s2}{\PYZdq{}}
\PY{n}{type}\PY{+w}{ }\PY{o}{=}\PY{+w}{ }\PY{l+s+s2}{\PYZdq{}}\PY{l+s+s2}{string}\PY{l+s+s2}{\PYZdq{}}
\PY{n}{escape}\PY{+w}{ }\PY{o}{=}\PY{+w}{ }\PY{l+s+s2}{\PYZdq{}}\PY{l+s+s2}{latex}\PY{l+s+s2}{\PYZdq{}}

\PY{k}{[}\PY{k}{latex}\PY{k}{.}\PY{k}{template}\PY{k}{.}\PY{k}{attributes}\PY{k}{.}\PY{k}{language}\PY{k}{]}
\PY{n}{default}\PY{+w}{ }\PY{o}{=}\PY{+w}{ }\PY{l+s+s2}{\PYZdq{}}\PY{l+s+s2}{english}\PY{l+s+s2}{\PYZdq{}}
\PY{n}{type}\PY{+w}{ }\PY{o}{=}\PY{+w}{ }\PY{l+s+s2}{\PYZdq{}}\PY{l+s+s2}{string}\PY{l+s+s2}{\PYZdq{}}
\PY{n}{normaliser}\PY{+w}{ }\PY{o}{=}\PY{+w}{ }\PY{l+s+s2}{\PYZdq{}}\PY{l+s+s2}{babel\PYZus{}language}\PY{l+s+s2}{\PYZdq{}}

\PY{k}{[}\PY{k}{latex}\PY{k}{.}\PY{k}{template}\PY{k}{.}\PY{k}{attributes}\PY{k}{.}\PY{k}{paper}\PY{k}{]}
\PY{n}{default}\PY{+w}{ }\PY{o}{=}\PY{+w}{ }\PY{l+s+s2}{\PYZdq{}}\PY{l+s+s2}{a4paper}\PY{l+s+s2}{\PYZdq{}}
\PY{n}{type}\PY{+w}{ }\PY{o}{=}\PY{+w}{ }\PY{l+s+s2}{\PYZdq{}}\PY{l+s+s2}{string}\PY{l+s+s2}{\PYZdq{}}
\PY{n}{normaliser}\PY{+w}{ }\PY{o}{=}\PY{+w}{ }\PY{l+s+s2}{\PYZdq{}}\PY{l+s+s2}{paper\PYZus{}option}\PY{l+s+s2}{\PYZdq{}}

\PY{k}{[}\PY{k}{latex}\PY{k}{.}\PY{k}{template}\PY{k}{.}\PY{k}{attributes}\PY{k}{.}\PY{k}{orientation}\PY{k}{]}
\PY{n}{default}\PY{+w}{ }\PY{o}{=}\PY{+w}{ }\PY{l+s+s2}{\PYZdq{}}\PY{l+s+s2}{portrait}\PY{l+s+s2}{\PYZdq{}}
\PY{n}{type}\PY{+w}{ }\PY{o}{=}\PY{+w}{ }\PY{l+s+s2}{\PYZdq{}}\PY{l+s+s2}{string}\PY{l+s+s2}{\PYZdq{}}
\PY{n}{normaliser}\PY{+w}{ }\PY{o}{=}\PY{+w}{ }\PY{l+s+s2}{\PYZdq{}}\PY{l+s+s2}{orientation}\PY{l+s+s2}{\PYZdq{}}
\PY{n}{choices}\PY{+w}{ }\PY{o}{=}\PY{+w}{ }\PY{p}{[}\PY{l+s+s2}{\PYZdq{}}\PY{l+s+s2}{portrait}\PY{l+s+s2}{\PYZdq{}}\PY{p}{,}\PY{+w}{ }\PY{l+s+s2}{\PYZdq{}}\PY{l+s+s2}{landscape}\PY{l+s+s2}{\PYZdq{}}\PY{p}{]}

\PY{k}{[}\PY{k}{latex}\PY{k}{.}\PY{k}{template}\PY{k}{.}\PY{k}{attributes}\PY{k}{.}\PY{k}{bibliography\PYZus{}style}\PY{k}{]}
\PY{n}{default}\PY{+w}{ }\PY{o}{=}\PY{+w}{ }\PY{l+s+s2}{\PYZdq{}}\PY{l+s+s2}{numeric}\PY{l+s+s2}{\PYZdq{}}
\PY{n}{type}\PY{+w}{ }\PY{o}{=}\PY{+w}{ }\PY{l+s+s2}{\PYZdq{}}\PY{l+s+s2}{string}\PY{l+s+s2}{\PYZdq{}}

\PY{c+c1}{\PYZsh{} Slots definition:}
\PY{c+c1}{\PYZsh{} \PYZhy{}\PYZhy{}\PYZhy{}\PYZhy{}\PYZhy{}\PYZhy{}\PYZhy{}\PYZhy{}\PYZhy{}\PYZhy{}\PYZhy{}\PYZhy{}\PYZhy{}\PYZhy{}\PYZhy{}\PYZhy{}\PYZhy{}}
\PY{c+c1}{\PYZsh{} Each slot represents a section of the document that can be filled}
\PY{c+c1}{\PYZsh{} with content. The \PYZsq{}depth\PYZsq{} parameter indicates the LaTeX sectioning}
\PY{c+c1}{\PYZsh{} command to use.}
\PY{k}{[}\PY{k}{latex}\PY{k}{.}\PY{k}{template}\PY{k}{.}\PY{k}{slots}\PY{k}{.}\PY{k}{mainmatter}\PY{k}{]}
\PY{n}{default}\PY{+w}{ }\PY{o}{=}\PY{+w}{ }\PY{k+kc}{true}
\PY{n}{depth}\PY{+w}{ }\PY{o}{=}\PY{+w}{ }\PY{l+s+s2}{\PYZdq{}}\PY{l+s+s2}{section}\PY{l+s+s2}{\PYZdq{}}

\PY{k}{[}\PY{k}{latex}\PY{k}{.}\PY{k}{template}\PY{k}{.}\PY{k}{slots}\PY{k}{.}\PY{k}{abstract}\PY{k}{]}
\PY{n}{depth}\PY{+w}{ }\PY{o}{=}\PY{+w}{ }\PY{l+s+s2}{\PYZdq{}}\PY{l+s+s2}{section}\PY{l+s+s2}{\PYZdq{}}
\PY{n}{strip\PYZus{}heading}\PY{+w}{ }\PY{o}{=}\PY{+w}{ }\PY{k+kc}{true}

\PY{c+c1}{\PYZsh{} Additional Assets:}
\PY{c+c1}{\PYZsh{} \PYZhy{}\PYZhy{}\PYZhy{}\PYZhy{}\PYZhy{}\PYZhy{}\PYZhy{}\PYZhy{}\PYZhy{}\PYZhy{}\PYZhy{}\PYZhy{}\PYZhy{}\PYZhy{}\PYZhy{}\PYZhy{}\PYZhy{}\PYZhy{}}
\PY{c+c1}{\PYZsh{} All additional assets to be included in the template package.}
\PY{c+c1}{\PYZsh{} They will be copied to the working directory when the template is used.}
\PY{k}{[}\PY{k}{latex}\PY{k}{.}\PY{k}{template}\PY{k}{.}\PY{k}{assets}\PY{k}{]}
\PY{l+s+s2}{\PYZdq{}}\PY{l+s+s2}{.latexmkrc}\PY{l+s+s2}{\PYZdq{}}\PY{+w}{ }\PY{o}{=}\PY{+w}{ }\PY{p}{\PYZob{}}\PY{+w}{ }\PY{n}{source}\PY{+w}{ }\PY{p}{=}\PY{+w}{ }\PY{l+s+s2}{\PYZdq{}}\PY{l+s+s2}{template/assets/latexmkrc}\PY{l+s+s2}{\PYZdq{}}\PY{+w}{ }\PY{p}{\PYZcb{}}

\PY{c+c1}{\PYZsh{} The Typst half, for a template that also serves \PYZhy{}\PYZhy{}format typst. Same shape:}
\PY{c+c1}{\PYZsh{} its own entrypoint, attributes and slots. A template without it still}
\PY{c+c1}{\PYZsh{} renders to Typst, in the standalone preamble.}
\PY{k}{[}\PY{k}{typst}\PY{k}{.}\PY{k}{template}\PY{k}{]}
\PY{n}{entrypoint}\PY{+w}{ }\PY{o}{=}\PY{+w}{ }\PY{l+s+s2}{\PYZdq{}}\PY{l+s+s2}{template/template.typ}\PY{l+s+s2}{\PYZdq{}}
\end{Verbatim}
\end{tscode}
\subsubsection{The Typst scaffolding's own variables}
Besides the attributes and the slots, a \tscodeinline{template.typ} receives:
\begin{itemize}
\item \tscodeinline{mainmatter} --- the \tscodeinline{texsmith.typ} contract definitions, the document's
callout-style choice, and the body of the default slot, in one piece. A
scaffolding with nothing to say about the contracts emits just this.
\item \tscodeinline{prelude} and \tscodeinline{body} --- the same two halves separately. A \tscodeinline{\#let ts-\allowbreak{}…}
redefinition only reaches the calls that follow it, so a template that
restyles a construct (\tscodeinline{texsmith.typ} says it may) writes its own definitions
between them:
\begin{tscode}[lang=jinja, engine=pygments]
\begin{Verbatim}[commandchars=\\\{\},breaklines, breakanywhere, commandchars=\\\{\}]
\PY{c+cp}{\PYZob{}\PYZob{}} \PY{n+nv}{prelude} \PY{c+cp}{\PYZcb{}\PYZcb{}}

\PY{x}{\PYZsh{}let ts\PYZhy{}code(..args) = block(stroke: 0.4pt, args.pos().at(0, default: []))}

\PY{c+cp}{\PYZob{}\PYZob{}} \PY{n+nv}{body} \PY{c+cp}{\PYZcb{}\PYZcb{}}
\end{Verbatim}
\end{tscode}
\end{itemize}
\subsubsection{Attribute schema}
Each attribute table accepts the following keys:
\begin{itemize}
\item \tscodeinline{default}: (required) the value injected when no override is provided.
\item \tscodeinline{type}: optional primitive among \tscodeinline{string}, \tscodeinline{integer}, \tscodeinline{float}, \tscodeinline{boolean}, \tscodeinline{list}, \tscodeinline{mapping}, \tscodeinline{any}. When omitted TeXSmith infers a type from \tscodeinline{default}.
\item \tscodeinline{sources}: the lookup paths (dot notation) used to pull overrides from document front matter. If omitted, TeXSmith probes \tscodeinline{press.press.<name>}, \tscodeinline{press.<name>}, then \tscodeinline{<name>}.
\item \tscodeinline{allow\_empty}: when \tscodeinline{false}, empty strings are coerced to \tscodeinline{None} so templates fall back to defaults.
\item \tscodeinline{required}: when \tscodeinline{true}, a missing override raises a \tscodeinline{TemplateError}.
\item \tscodeinline{choices}: restricts string/integer values to the provided list.
\item \tscodeinline{escape}: currently supports \tscodeinline{latex} to automatically escape special characters coming from user input.
\item \tscodeinline{format}: what the value \emph{is}. \tscodeinline{markdown} (the default) renders the value --- a line or two of prose --- through the writer; \tscodeinline{raw} emits it as it stands; \tscodeinline{file} reads the value as the \textbf{path of a file} and inlines its text verbatim. A \tscodeinline{file} attribute cannot also be escaped, its default is the empty string, and a value that is not a string (a mapping the template reads through other attributes) leaves it at that default. Relative paths start at the converted document's directory, and at the project directory for a book built from a site's configuration file --- where \tscodeinline{copy\_files} starts too. A path that names no file is a \tscodeinline{TemplateError}, so a missing title page never silently falls back to the template's own.
\item \tscodeinline{normaliser}: a built-in normaliser name \textbf{or} a \tscodeinline{"module:callable"} import reference pointing at your own callable that post-processes the value (see below).
\end{itemize}
Attributes are resolved before \tscodeinline{WrappableTemplate.prepare\_context} runs, meaning template classes only need to handle presentation-specific tweaks (for example, combining authors, building option strings, or removing temporary keys).
\subsection{Attribute ownership \& precedence}
\begin{itemize}
\item Each attribute belongs to a single owner (template by default, fragment when declared in \tscodeinline{fragment.toml}). Conflicting owners raise a \tscodeinline{TemplateError}.
\item Overrides flow: CLI/front matter \tsscript{symbols}{\(\rightarrow\) }attribute resolver (type coercion, normalisers) \tsscript{symbols}{\(\rightarrow\) }template emitters/fragment defaults. Empty strings are dropped when \tscodeinline{allow\_empty = false}.
\item Attributes no longer live under ad-hoc dotted names (\tscodeinline{press.*}); normalisation collapses supported aliases onto the declared attribute name.
\end{itemize}
\subsection{Template slots and built-ins}
\begin{itemize}
\item Templates must declare at least one slot; \tscodeinline{mainmatter} is the default when absent. Slots surface \tscodeinline{depth}, \tscodeinline{offset}, \tscodeinline{strip\_heading}, and optional \tscodeinline{base\_level}.
\item Built-in \tscodeinline{article} exposes \tscodeinline{mainmatter}, \tscodeinline{abstract}, \tscodeinline{appendix}, and \tscodeinline{backmatter}. \tscodeinline{book} inserts \tscodeinline{appendix} before \tscodeinline{backmatter} and supports \tscodeinline{part} toggling. \tscodeinline{letter} uses \tscodeinline{mainmatter} only.
\item Deprecated attributes removed from built-ins: \tscodeinline{article}/\tscodeinline{book} no longer accept \tscodeinline{cover}, \tscodeinline{covercolor}, \tscodeinline{twocolumn}, or template-owned \tscodeinline{emoji}. Backmatter/preamble overrides moved into fragments where applicable.
\end{itemize}
\subsubsection{Built-in normalisers}
\begin{center}
\begin{tabularx}{\linewidth}{>{\raggedright\arraybackslash}X>{\raggedright\arraybackslash}X}
\toprule
\textbf{Name} & \textbf{Purpose} \\
\midrule
\tscodeinline{paper\_option} & Validates paper sizes and emits \tscodeinline{<size>paper} strings (for example \tscodeinline{letterpaper}). \\
\tscodeinline{orientation} & Normalises orientation flags and guarantees \tscodeinline{portrait} or \tscodeinline{landscape}. \\
\tscodeinline{babel\_language} & Maps ISO-like language codes to Babel identifiers (for example \tscodeinline{fr} \tsscript{symbols}{\(\rightarrow\) }\tscodeinline{french}). \\
\bottomrule
\end{tabularx}
\end{center}
Normalisers run after type coercion and before escaping. They can also reuse existing defaults to provide fallbacks (\tscodeinline{paper\_option} and \tscodeinline{orientation} do this).
\subsubsection{Custom normalisers from template packages}
Beyond the built-in names, \tscodeinline{normaliser} accepts a fully-qualified callable
reference so a template package can supply its own normaliser without touching
TeXSmith internals. Use the \tscodeinline{"package.module:callable"} form (the dotted
\tscodeinline{"package.module.callable"} form is also accepted):
\begin{tscode}[lang=toml, engine=pygments]
\begin{Verbatim}[commandchars=\\\{\},breaklines, breakanywhere, commandchars=\\\{\}]
\PY{k}{[}\PY{k}{latex}\PY{k}{.}\PY{k}{template}\PY{k}{.}\PY{k}{attributes}\PY{k}{.}\PY{k}{rules}\PY{k}{]}
\PY{n}{type}\PY{+w}{ }\PY{o}{=}\PY{+w}{ }\PY{l+s+s2}{\PYZdq{}}\PY{l+s+s2}{any}\PY{l+s+s2}{\PYZdq{}}
\PY{n}{normaliser}\PY{+w}{ }\PY{o}{=}\PY{+w}{ }\PY{l+s+s2}{\PYZdq{}}\PY{l+s+s2}{texsmith\PYZus{}template\PYZus{}exam.rules:resolve\PYZus{}attribute}\PY{l+s+s2}{\PYZdq{}}
\end{Verbatim}
\end{tscode}
The referenced object is imported (and cached) when the manifest is loaded, and
must honour the normaliser contract:
\begin{tscode}[lang=python, engine=pygments]
\begin{Verbatim}[commandchars=\\\{\},breaklines, breakanywhere, commandchars=\\\{\}]
\PY{k}{def}\PY{+w}{ }\PY{n+nf}{resolve\PYZus{}attribute}\PY{p}{(}\PY{n}{value}\PY{p}{,} \PY{n}{spec}\PY{p}{,} \PY{n}{fallback}\PY{p}{)}\PY{p}{:}
    \PY{c+c1}{\PYZsh{} value:    the already type\PYZhy{}coerced attribute value}
    \PY{c+c1}{\PYZsh{} spec:     the TemplateAttributeSpec (use spec.name in error messages)}
    \PY{c+c1}{\PYZsh{} fallback: the attribute\PYZsq{}s resolved default}
    \PY{c+c1}{\PYZsh{} returns:  the normalised value}
    \PY{o}{.}\PY{o}{.}\PY{o}{.}
\end{Verbatim}
\end{tscode}
If the module or attribute cannot be found, or the target is not callable,
TeXSmith raises a \tscodeinline{TemplateError} naming the offending template, attribute, and
\tscodeinline{normaliser} value.

Authors who prefer a short registry name can instead call
\tscodeinline{texsmith.core.templates.register\_attribute\_normaliser(name, func)} at import
time and reference it by \tscodeinline{name}. Registering over a built-in name raises a
\tscodeinline{TemplateError} unless \tscodeinline{override=True} is passed. The \tscodeinline{"module:callable"} form is
collision-free and needs no registration, so it is generally preferred.
\subsubsection{Metadata overrides}
At runtime TeXSmith merges template defaults with document front matter. By default, overrides come from the \tscodeinline{press} section:
\begin{tscode}[lang=yaml, engine=pygments]
\begin{Verbatim}[commandchars=\\\{\},breaklines, breakanywhere, commandchars=\\\{\}]
\PY{n+nn}{\PYZhy{}\PYZhy{}\PYZhy{}}
\PY{n+nt}{press}\PY{p}{:}
\PY{+w}{  }\PY{n+nt}{title}\PY{p}{:}\PY{+w}{ }\PY{l+lScalar+lScalarPlain}{Sample}\PY{l+lScalar+lScalarPlain}{ }\PY{l+lScalar+lScalarPlain}{Article}
\PY{+w}{  }\PY{n+nt}{subtitle}\PY{p}{:}\PY{+w}{ }\PY{l+lScalar+lScalarPlain}{Insights}\PY{l+lScalar+lScalarPlain}{ }\PY{l+lScalar+lScalarPlain}{on}\PY{l+lScalar+lScalarPlain}{ }\PY{l+lScalar+lScalarPlain}{Cheese}
\PY{+w}{  }\PY{n+nt}{language}\PY{p}{:}\PY{+w}{ }\PY{l+lScalar+lScalarPlain}{fr}
\PY{+w}{  }\PY{n+nt}{authors}\PY{p}{:}
\PY{+w}{    }\PY{p+pIndicator}{\PYZhy{}}\PY{+w}{ }\PY{n+nt}{name}\PY{p}{:}\PY{+w}{ }\PY{l+lScalar+lScalarPlain}{Ada}\PY{l+lScalar+lScalarPlain}{ }\PY{l+lScalar+lScalarPlain}{Lovelace}
\PY{+w}{      }\PY{n+nt}{affiliation}\PY{p}{:}\PY{+w}{ }\PY{l+lScalar+lScalarPlain}{Analytical}\PY{l+lScalar+lScalarPlain}{ }\PY{l+lScalar+lScalarPlain}{Engine}
\PY{+w}{    }\PY{p+pIndicator}{\PYZhy{}}\PY{+w}{ }\PY{n+nt}{name}\PY{p}{:}\PY{+w}{ }\PY{l+lScalar+lScalarPlain}{Grace}\PY{l+lScalar+lScalarPlain}{ }\PY{l+lScalar+lScalarPlain}{Hopper}
\PY{n+nn}{\PYZhy{}\PYZhy{}\PYZhy{}}
\end{Verbatim}
\end{tscode}
With the manifest above TeXSmith will:
\begin{itemize}
\item Escape \tscodeinline{title} and \tscodeinline{subtitle} for \LaTeX{} compatibility.
\item Map \tscodeinline{language: fr} to \tscodeinline{french}.
\item Normalise author collections so individual templates can format them consistently.
\item Populate both defaults (\tscodeinline{author}) and structured data (\tscodeinline{authors}) that template classes can consume.
\end{itemize}
To pull metadata from alternative locations add explicit \tscodeinline{sources}:
\begin{tscode}[lang=toml, engine=pygments]
\begin{Verbatim}[commandchars=\\\{\},breaklines, breakanywhere, commandchars=\\\{\}]
\PY{k}{[}\PY{k}{latex}\PY{k}{.}\PY{k}{template}\PY{k}{.}\PY{k}{attributes}\PY{k}{.}\PY{k}{project\PYZus{}code}\PY{k}{]}
\PY{n}{default}\PY{+w}{ }\PY{o}{=}\PY{+w}{ }\PY{l+s+s2}{\PYZdq{}}\PY{l+s+s2}{\PYZdq{}}
\PY{n}{type}\PY{+w}{ }\PY{o}{=}\PY{+w}{ }\PY{l+s+s2}{\PYZdq{}}\PY{l+s+s2}{string}\PY{l+s+s2}{\PYZdq{}}
\PY{n}{sources}\PY{+w}{ }\PY{o}{=}\PY{+w}{ }\PY{p}{[}\PY{l+s+s2}{\PYZdq{}}\PY{l+s+s2}{press.project.code}\PY{l+s+s2}{\PYZdq{}}\PY{p}{,}\PY{+w}{ }\PY{l+s+s2}{\PYZdq{}}\PY{l+s+s2}{project.code}\PY{l+s+s2}{\PYZdq{}}\PY{p}{]}
\PY{n}{allow\PYZus{}empty}\PY{+w}{ }\PY{o}{=}\PY{+w}{ }\PY{k+kc}{false}
\end{Verbatim}
\end{tscode}
\subsubsection{Slots}
Slots control where converted content is injected. Each entry allows:
\begin{itemize}
\item \tscodeinline{depth}: maps to a \LaTeX{} sectioning command (\tscodeinline{section}, \tscodeinline{chapter}, etc.).
\item \tscodeinline{base\_level} and \tscodeinline{offset}: fine tune heading levels when rendering.
\item \tscodeinline{default}: mark exactly one slot as the primary sink for content.
\item \tscodeinline{strip\_heading}: remove the first heading when populating the slot (useful for abstracts).
\end{itemize}
Slots become Jinja variables inside the template (\tscodeinline{\textbackslash{}VAR\{abstract\}}, \tscodeinline{\textbackslash{}VAR\{mainmatter\}}).
\subsection{Overrides}
The \LaTeX{} body comes from tmark's writer, which emits a fixed macro or
environment per construct. Structural constructs (emphasis, lists, headings,
figures, tables, links, footnotes) are plain \LaTeX{} the writer owns; everything a
template may want to restyle is a \textbf{contract macro} provided by a \tscodeinline{ts-\allowbreak{}*}
fragment. Redefine it in your template's \tscodeinline{.tex}, after \tscodeinline{\textbackslash{}VAR\{extra\_packages\}}
(that is where the fragments are \tscodeinline{\textbackslash{}usepackage}d):
\begin{tscode}[lang=latex, engine=pygments]
\begin{Verbatim}[commandchars=\\\{\},breaklines, breakanywhere, commandchars=\\\{\}]
\PY{k}{\PYZbs{}VAR}\PY{n+nb}{\PYZob{}}extra\PY{n+nb}{\PYZus{}}packages\PY{n+nb}{\PYZcb{}}
\PY{k}{\PYZbs{}tcbset}\PY{n+nb}{\PYZob{}}/ts/code/.append style=\PY{n+nb}{\PYZob{}}frame hidden, boxrule=0pt\PY{n+nb}{\PYZcb{}}\PY{n+nb}{\PYZcb{}}
\PY{k}{\PYZbs{}RenewDocumentCommand}\PY{n+nb}{\PYZob{}}\PY{k}{\PYZbs{}tscodeinline}\PY{n+nb}{\PYZcb{}}\PY{n+nb}{\PYZob{}}O\PY{n+nb}{\PYZob{}}\PY{n+nb}{\PYZcb{}}m\PY{n+nb}{\PYZcb{}}\PY{n+nb}{\PYZob{}}\PY{k}{\PYZbs{}texttt}\PY{n+nb}{\PYZob{}}\PYZsh{}2\PY{n+nb}{\PYZcb{}}\PY{n+nb}{\PYZcb{}}
\PY{k}{\PYZbs{}RenewDocumentEnvironment}\PY{n+nb}{\PYZob{}}tsdiv@multicolumn\PY{n+nb}{\PYZcb{}}\PY{n+nb}{\PYZob{}}O\PY{n+nb}{\PYZob{}}\PY{n+nb}{\PYZcb{}}\PY{n+nb}{\PYZcb{}}\PY{n+nb}{\PYZob{}}\PY{k}{\PYZbs{}begin}\PY{n+nb}{\PYZob{}}multicols\PY{n+nb}{\PYZcb{}}\PY{n+nb}{\PYZob{}}3\PY{n+nb}{\PYZcb{}}\PY{n+nb}{\PYZcb{}}\PY{n+nb}{\PYZob{}}\PY{k}{\PYZbs{}end}\PY{n+nb}{\PYZob{}}multicols\PY{n+nb}{\PYZcb{}}\PY{n+nb}{\PYZcb{}}
\end{Verbatim}
\end{tscode}
See Contract macros for the full table of macros, their keys and
the three levels of override.
\begin{tscallout}[kind=warning, title={Jinja partials are no longer the rendering layer}]
\tscodeinline{latex.template.override}, fragment \tscodeinline{partials} and \tscodeinline{required\_partials} were
deprecated in 0.7.0 and are removed in 0.8.0: the manifest no longer reads
them. \href{partials.md\#former-partials}{Contract macros} maps every former
partial to its replacement macro.
\end{tscallout}
\subsection{A cover, an imprint and a preamble of your own}
The \tscodeinline{book} template takes three attributes so a project can keep its own front
matter without forking the template:
\begin{center}
\begin{tabularx}{\linewidth}{>{\raggedright\arraybackslash}X>{\raggedright\arraybackslash}X>{\raggedright\arraybackslash}X}
\toprule
\textbf{Attribute} & \textbf{Value} & \textbf{Where it lands} \\
\midrule
\tscodeinline{press.preamble} & \LaTeX{}, inline & the end of the preamble, after the fragments --- load a package, redefine a contract macro, override anything above \\
\tscodeinline{press.titlepage} & a path & in place of \tscodeinline{\textbackslash{}maketitle}, inside \tscodeinline{\textbackslash{}frontmatter} \\
\tscodeinline{press.imprint} & a path & in place of the imprint page built from \tscodeinline{press.imprint.thanks} / \tscodeinline{.license} / \tscodeinline{.copyright} \\
\bottomrule
\end{tabularx}
\end{center}
\tscodeinline{press.imprint} is therefore either a path --- your own page --- or the mapping the
template renders itself; it is never both.
\begin{tscode}[lang=yaml, engine=pygments]
\begin{Verbatim}[commandchars=\\\{\},breaklines, breakanywhere, commandchars=\\\{\}]
\PY{n+nn}{\PYZhy{}\PYZhy{}\PYZhy{}}
\PY{n+nt}{press}\PY{p}{:}
\PY{+w}{  }\PY{n+nt}{titlepage}\PY{p}{:}\PY{+w}{ }\PY{l+lScalar+lScalarPlain}{tex/titlepage.tex}
\PY{+w}{  }\PY{n+nt}{imprint}\PY{p}{:}\PY{+w}{ }\PY{l+lScalar+lScalarPlain}{tex/imprint.tex}
\PY{+w}{  }\PY{n+nt}{preamble}\PY{p}{:}\PY{+w}{ }\PY{p+pIndicator}{|}
\PY{+w}{    }\PY{n+no}{\PYZbs{}usepackage\PYZob{}acmelogo\PYZcb{}}
\PY{n+nn}{\PYZhy{}\PYZhy{}\PYZhy{}}
\end{Verbatim}
\end{tscode}
The files are the author's \LaTeX{} and are inlined verbatim, so they can use the
macros the template defines --- \tscodeinline{\textbackslash{}booktitle}, \tscodeinline{\textbackslash{}booksubtitle}, \tscodeinline{\textbackslash{}bookauthor},
\tscodeinline{\textbackslash{}bookemail}, \tscodeinline{\textbackslash{}bookpublisher}, \tscodeinline{\textbackslash{}bookedition}, \tscodeinline{\textbackslash{}bookdate} --- and anything a
\tscodeinline{.sty} of yours provides:
\begin{tscode}[lang=latex, engine=pygments]
\begin{Verbatim}[commandchars=\\\{\},breaklines, breakanywhere, commandchars=\\\{\}]
\PY{k}{\PYZbs{}begin}\PY{n+nb}{\PYZob{}}titlingpage\PY{n+nb}{\PYZcb{}}
  \PY{k}{\PYZbs{}acmelogo}
  \PY{k}{\PYZbs{}begin}\PY{n+nb}{\PYZob{}}center\PY{n+nb}{\PYZcb{}}
    \PY{n+nb}{\PYZob{}}\PY{k}{\PYZbs{}huge}\PY{k}{\PYZbs{}bfseries} \PY{k}{\PYZbs{}booktitle}\PY{n+nb}{\PYZcb{}}\PY{k}{\PYZbs{}\PYZbs{}}\PY{n+na}{[1.5cm]}
    \PY{n+nb}{\PYZob{}}\PY{k}{\PYZbs{}Large} \PY{k}{\PYZbs{}bookauthor}\PY{n+nb}{\PYZcb{}}\PY{k}{\PYZbs{}\PYZbs{}}\PY{n+na}{[5pt]}
    \PY{k}{\PYZbs{}bookemail}\PY{k}{\PYZbs{}\PYZbs{}}
    \PY{k}{\PYZbs{}vfill}
    \PY{k}{\PYZbs{}bookdate}
  \PY{k}{\PYZbs{}end}\PY{n+nb}{\PYZob{}}center\PY{n+nb}{\PYZcb{}}
\PY{k}{\PYZbs{}end}\PY{n+nb}{\PYZob{}}titlingpage\PY{n+nb}{\PYZcb{}}
\end{Verbatim}
\end{tscode}
A \tscodeinline{.sty} the preamble loads has to reach the bundle: ship it as a template
asset (\tscodeinline{[latex.template.assets]}), or, for a book built from a site, list it
under the book's \tscodeinline{copy\_files}, which copies it next to the generated \tscodeinline{.tex}
where the engine looks:
\begin{tscode}[lang=yaml, engine=pygments]
\begin{Verbatim}[commandchars=\\\{\},breaklines, breakanywhere, commandchars=\\\{\}]
\PY{n+nt}{plugins}\PY{p}{:}
\PY{+w}{  }\PY{p+pIndicator}{\PYZhy{}}\PY{+w}{ }\PY{n+nt}{texsmith}\PY{p}{:}
\PY{+w}{      }\PY{n+nt}{books}\PY{p}{:}
\PY{+w}{        }\PY{p+pIndicator}{\PYZhy{}}\PY{+w}{ }\PY{n+nt}{title}\PY{p}{:}\PY{+w}{ }\PY{l+lScalar+lScalarPlain}{The}\PY{l+lScalar+lScalarPlain}{ }\PY{l+lScalar+lScalarPlain}{Book}
\PY{+w}{          }\PY{n+nt}{press}\PY{p}{:}
\PY{+w}{            }\PY{n+nt}{titlepage}\PY{p}{:}\PY{+w}{ }\PY{l+lScalar+lScalarPlain}{tex/titlepage.tex}
\PY{+w}{            }\PY{n+nt}{imprint}\PY{p}{:}\PY{+w}{ }\PY{l+lScalar+lScalarPlain}{tex/imprint.tex}
\PY{+w}{            }\PY{n+nt}{preamble}\PY{p}{:}\PY{+w}{ }\PY{l+lScalar+lScalarPlain}{\PYZbs{}usepackage\PYZob{}acmelogo\PYZcb{}}
\PY{+w}{          }\PY{n+nt}{copy\PYZus{}files}\PY{p}{:}
\PY{+w}{            }\PY{n+nt}{tex/*.sty}\PY{p}{:}\PY{+w}{ }\PY{l+lScalar+lScalarPlain}{.}
\end{Verbatim}
\end{tscode}
\subsection{Slot strategies}
Slots determine how multiple Markdown documents (or sections) flow into the
\LaTeX{} structure. Typical patterns:
\begin{itemize}
\item \textbf{Single document, single slot} -- map the only input with \tscodeinline{-\allowbreak{}-\allowbreak{}slot mainmatter:@document}.
\item \textbf{Front matter + main matter} -- convert two files (eg \tscodeinline{intro.md},
\tscodeinline{book.md}) and pass \tscodeinline{-\allowbreak{}-\allowbreak{}slot frontmatter:intro.md} and \tscodeinline{-\allowbreak{}-\allowbreak{}slot mainmatter:book.md}.
\item \textbf{Selective sections} -- reference headings or IDs:
\tscodeinline{-\allowbreak{}-\allowbreak{}slot abstract:paper.md:@abstract -\allowbreak{}-\allowbreak{}slot mainmatter:paper.md:"Results"} injects
only the abstract heading and the ``Results'' section into separate slots.
\item \textbf{Per-language appendices} -- define slots (\tscodeinline{appendix\_en}, \tscodeinline{appendix\_fr}) and use front-matter metadata (\tscodeinline{press.slot.appendix\_en: docs/en.md}) so automation scripts do not need to pass CLI flags.
\end{itemize}
Slots are resolved by the document slot mapping across CLI flags, front matter (\tscodeinline{press.slot.*}), and API overrides, so mix and match whichever suits your workflow.
\subsection{Testing your template}
Before publishing, run the built-in test suite against your template:
\begin{tscode}[lang=bash, engine=pygments]
\begin{Verbatim}[commandchars=\\\{\},breaklines, breakanywhere, commandchars=\\\{\}]
uv\PY{+w}{ }run\PY{+w}{ }pytest\PY{+w}{ }tests/test\PYZus{}template\PYZus{}attributes.py
\end{Verbatim}
\end{tscode}
This repository includes examples for the bundled templates; copy one into your package and adjust expectations to match your manifest. Automated tests help ensure attributes, slots, and metadata mappings keep working as the TeXSmith runtime evolves.
\section{Next steps}
\begin{itemize}
\item Study the Template Cookbook for practical recipes (title pages, metadata bindings, bibliography tweaks).
\item Browse the API high-level guide to orchestrate templates programmatically with \tscodeinline{ConversionService}.
\end{itemize}
\end{document}
